\section{Tổng quan chức năng}
\label{section:2.2}

Dựa trên phạm vi đã xác định, GoPortal được xây dựng như một nền tảng cộng đồng game trực tuyến, hỗ trợ người dùng quản lý tài khoản, kết nối bạn bè, tham gia máy chủ, trao đổi tin nhắn, chơi game, đồng bộ trạng thái và tham gia giải đấu. Hệ thống có nhiều nhóm tác nhân khác nhau như khách, người dùng, chủ máy chủ, nhà phát triển và quản trị viên. Mỗi tác nhân được gắn với những nhóm chức năng phù hợp với vai trò của mình trong hệ thống.

Phần này trình bày tổng quan các chức năng thông qua biểu đồ use case tổng quát và các biểu đồ use case phân rã. Các nội dung ở đây chỉ mô tả chức năng mức cao, chưa đi vào luồng xử lý chi tiết. Phần đặc tả chi tiết cho từng use case quan trọng sẽ được trình bày ở phần \ref{section:2.3}.

\subsection{Biểu đồ use case tổng quát}
\label{subsection:2.2.1}

Biểu đồ use case tổng quát dưới đây thể hiện các tác nhân chính và các nhóm chức năng lớn của hệ thống GoPortal.

\begin{figure}[H]
    \centering
    \includegraphics[width=\textwidth]{Hinhve/Chuong2/Usecase/Tong-quan.png}
    \caption{Biểu đồ use case tổng quát}
    \label{fig:2-1-usecase-tong-quat}
\end{figure}

Hình \ref{fig:2-1-usecase-tong-quat} xác định năm tác nhân chính gồm Khách, Người dùng, Chủ máy chủ, Nhà phát triển và Quản trị viên. Trong đó, Người dùng là tác nhân trung tâm của hệ thống. Chủ máy chủ và Nhà phát triển là hai vai trò được mở rộng từ Người dùng, nghĩa là ngoài các chức năng cơ bản của Người dùng, các tác nhân này còn có thêm những chức năng chuyên biệt.

Khách là người chưa đăng nhập hoặc chưa có tài khoản, chỉ có thể sử dụng một số chức năng công khai như Chơi game và Đồng bộ trong phạm vi được hệ thống cho phép. Người dùng sau khi đăng nhập có thể Quản lý thông tin cá nhân, Quản lý bạn bè, Tham gia máy chủ, Quản lý tin nhắn, Tham gia giải đấu và Chơi game, Đồng bộ. Đây là các chức năng cơ bản phục vụ nhu cầu tương tác, giao tiếp và tham gia hoạt động game trong cộng đồng.

Chủ máy chủ có thêm các chức năng Quản lý máy chủ và Thành viên, Quản lý kênh và Phân quyền, đồng thời có thể Quản lý giải đấu trong phạm vi máy chủ. Nhà phát triển thực hiện chức năng Phát triển và Đăng tải game, phục vụ nhu cầu đưa game lên hệ thống để người dùng khác trải nghiệm. Quản trị viên đảm nhiệm nhóm chức năng Quản trị người dùng và Chợ game, nhằm kiểm soát người dùng, nội dung game và các hoạt động quản trị ở cấp toàn hệ thống.

Nhìn chung, biểu đồ tổng quát cho thấy GoPortal có bốn nhóm nghiệp vụ chính: quản lý tài khoản và quan hệ xã hội, quản lý cộng đồng máy chủ, chơi và đăng tải game, quản lý giải đấu và chợ game. Các nhóm chức năng này là cơ sở để tiếp tục phân rã và đặc tả trong các mục sau.

\subsection{Biểu đồ use case phân rã Quản lý tài khoản}
\label{subsection:2.2.2}

Use case Quản lý tài khoản mô tả các chức năng liên quan đến việc tạo tài khoản, đăng nhập, đăng xuất và quản lý hồ sơ cá nhân.

\begin{figure}[H]
    \centering
    \includegraphics[width=\textwidth]{Hinhve/Chuong2/Usecase/Quan-ly-tai-khoan.png}
    \caption{Biểu đồ use case phân rã Quản lý tài khoản}
    \label{fig:2-2-usecase-quan-ly-tai-khoan}
\end{figure}

Hình \ref{fig:2-2-usecase-quan-ly-tai-khoan} gồm hai tác nhân là Khách và Người dùng. Khách có thể Đăng ký tài khoản để tạo tài khoản mới hoặc Đăng nhập hệ thống để xác thực và chuyển sang vai trò Người dùng. Đây là hai chức năng đầu vào giúp người dùng bắt đầu sử dụng các chức năng cần xác thực của GoPortal.

Sau khi đăng nhập, Người dùng có thể Xem hồ sơ cá nhân, Cập nhật thông tin cá nhân, Đăng xuất và Đổi mật khẩu. Các chức năng này giúp người dùng kiểm tra thông tin hiển thị, cập nhật hồ sơ, kết thúc phiên làm việc khi cần và đảm bảo an toàn cho tài khoản. Nhóm chức năng Quản lý tài khoản là nền tảng cho các chức năng khác như tham gia máy chủ, gửi tin nhắn, kết bạn, đăng tải game và tham gia giải đấu.

\subsection{Biểu đồ use case phân rã Quản lý bạn bè}
\label{subsection:2.2.3}

Use case Quản lý bạn bè mô tả các chức năng giúp Người dùng thiết lập và kiểm soát quan hệ xã hội với các tài khoản khác trong hệ thống.

\begin{figure}[H]
    \centering
    \includegraphics[width=0.85\textwidth]{Hinhve/Chuong2/Usecase/Quan-ly-ban-be.png}
    \caption{Biểu đồ use case phân rã Quản lý bạn bè}
    \label{fig:2-3-usecase-quan-ly-ban-be}
\end{figure}

Hình \ref{fig:2-3-usecase-quan-ly-ban-be} cho thấy Người dùng là tác nhân thực hiện toàn bộ các chức năng trong nhóm Quản lý bạn bè. Người dùng có thể Xem danh sách bạn bè để theo dõi các tài khoản đã kết nối, Gửi yêu cầu kết bạn để thiết lập quan hệ mới, Chấp nhận kết bạn hoặc Từ chối kết bạn khi nhận được yêu cầu từ người khác.

Ngoài ra, Người dùng có thể Hủy kết bạn khi không muốn duy trì quan hệ đã có. Trong trường hợp muốn giới hạn tương tác với một tài khoản, Người dùng có thể Chặn người dùng và sau đó Bỏ chặn nếu muốn khôi phục khả năng tương tác. Nhóm chức năng này giúp người dùng chủ động quản lý mạng lưới quan hệ xã hội, tạo nền tảng cho các hoạt động cộng đồng như nhắn tin, tham gia máy chủ, chơi game cùng nhau và tham gia giải đấu.

\subsection{Biểu đồ use case phân rã Quản lý máy chủ}
\label{subsection:2.2.4}

Use case Quản lý máy chủ mô tả các chức năng liên quan đến việc người dùng tham gia máy chủ và chủ máy chủ quản lý không gian cộng đồng của mình.

\begin{figure}[H]
    \centering
    \includegraphics[width=\textwidth]{Hinhve/Chuong2/Usecase/Quan-ly-may-chu.png}
    \caption{Biểu đồ use case phân rã Quản lý máy chủ}
    \label{fig:2-4-usecase-quan-ly-may-chu}
\end{figure}

Hình \ref{fig:2-4-usecase-quan-ly-may-chu} cho thấy Người dùng có thể Xem danh sách máy chủ, Tham gia máy chủ, Gửi yêu cầu tham gia và Rời máy chủ. Các chức năng này giúp người dùng tìm kiếm, gia nhập hoặc rời khỏi các cộng đồng game phù hợp với nhu cầu của mình.

Đối với Chủ máy chủ, hệ thống cung cấp các chức năng Tạo máy chủ, Cập nhật máy chủ, Duyệt yêu cầu tham gia, Quản lý thành viên, Quản lý vai trò, Tạo kênh, Cập nhật kênh và Phân quyền kênh. Nhóm chức năng này giúp chủ máy chủ tổ chức cộng đồng, kiểm soát thành viên và thiết lập phạm vi truy cập cho từng kênh trong máy chủ.

\subsection{Biểu đồ use case phân rã Quản lý tin nhắn}
\label{subsection:2.2.5}

Use case Quản lý tin nhắn mô tả các chức năng phục vụ hoạt động trao đổi thông tin giữa người dùng trong các kênh của máy chủ.

\begin{figure}[H]
    \centering
    \includegraphics[width=\textwidth]{Hinhve/Chuong2/Usecase/Quan-ly-tin-nhan.png}
    \caption{Biểu đồ use case phân rã Quản lý tin nhắn}
    \label{fig:2-5-usecase-quan-ly-tin-nhan}
\end{figure}

Hình \ref{fig:2-5-usecase-quan-ly-tin-nhan} thể hiện tác nhân Người dùng với các chức năng Xem tin nhắn và Gửi tin nhắn. Đây là hai chức năng cơ bản nhất để người dùng theo dõi nội dung trao đổi và tham gia thảo luận trong cộng đồng.

Bên cạnh đó, Người dùng có thể Chỉnh sửa tin nhắn, Xóa tin nhắn, Trả lời tin nhắn, Thả phản ứng và Đánh dấu đã đọc. Các chức năng bổ sung này giúp quá trình giao tiếp linh hoạt hơn, đồng thời hỗ trợ người dùng quản lý trạng thái đọc và tương tác với nội dung trong kênh.

\subsection{Biểu đồ use case phân rã Chơi game và Đăng tải game}
\label{subsection:2.2.6}

Use case Chơi game và Đăng tải game mô tả nhóm chức năng liên quan đến việc khám phá, trải nghiệm game và đăng tải game lên hệ thống.

\begin{figure}[H]
    \centering
    \includegraphics[width=\textwidth]{Hinhve/Chuong2/Usecase/Choi-game-va-dang-tai-game.png}
    \caption{Biểu đồ use case phân rã Chơi game và Đăng tải game}
    \label{fig:2-6-usecase-choi-game-va-dang-tai-game}
\end{figure}

Hình \ref{fig:2-6-usecase-choi-game-va-dang-tai-game} cho thấy Khách có thể Xem danh sách game, Xem chi tiết game và Chơi game công khai trong phạm vi được phép. Người dùng sau khi đăng nhập có thể Chơi game, Tạo phiên chơi, Đồng bộ trạng thái, Chia sẻ điểm số và Chia sẻ thành tích. Các chức năng này giúp hoạt động chơi game gắn với tài khoản và tương tác cộng đồng.

Đối với Nhà phát triển, hệ thống hỗ trợ Tạo thông tin game, Cập nhật thông tin game, Tải lên bản build, Gửi kiểm duyệt và Quản lý game đã đăng. Nhóm chức năng này cho phép nhà phát triển đưa game lên GoPortal, cập nhật nội dung và chuẩn bị game để người dùng khác có thể trải nghiệm.

\subsection{Biểu đồ use case phân rã Quản lý và Tham gia giải đấu}
\label{subsection:2.2.7}

Use case Quản lý và Tham gia giải đấu được tách thành hai biểu đồ để giảm độ phức tạp: nhóm chức năng thiết lập, tham gia giải đấu và nhóm chức năng vận hành trận đấu trong giải đấu.

\begin{figure}[H]
    \centering
    \includegraphics[width=\textwidth]{Hinhve/Chuong2/Usecase/Thiet-lap-va-tham-gia-giai-dau.png}
    \caption{Biểu đồ use case phân rã Thiết lập và Tham gia giải đấu}
    \label{fig:2-7a-usecase-thiet-lap-va-tham-gia-giai-dau}
\end{figure}

Hình \ref{fig:2-7a-usecase-thiet-lap-va-tham-gia-giai-dau} mô tả giai đoạn thiết lập và tham gia giải đấu. Người tổ chức có thể tạo giải đấu, cập nhật thông tin, mở đăng ký, chuyển sang check-in, bắt đầu giải đấu, tạo bracket và kết thúc hoặc hủy giải đấu. Quản trị giải đấu hỗ trợ gán vai trò và quản lý người tham gia trong phạm vi giải đấu.

Người tham gia có thể xem giải đấu, đăng ký cá nhân, hủy đăng ký và check-in khi giải đấu bước vào giai đoạn tương ứng. Với giải đấu theo đội, Đội trưởng có thêm các chức năng tạo đội, quản lý thành viên đội và check-in cho đội.

\begin{figure}[H]
    \centering
    \includegraphics[width=\textwidth]{Hinhve/Chuong2/Usecase/Van-hanh-tran-dau-giai-dau.png}
    \caption{Biểu đồ use case phân rã Vận hành trận đấu trong giải đấu}
    \label{fig:2-7b-usecase-van-hanh-tran-dau-giai-dau}
\end{figure}

Hình \ref{fig:2-7b-usecase-van-hanh-tran-dau-giai-dau} mô tả các chức năng sau khi giải đấu đã bắt đầu. Người tổ chức hoặc Quản trị giải đấu có thể bắt đầu trận đấu, cập nhật trạng thái trận, tạo workspace trận và ghi đè kết quả khi cần. Người tham gia và Đội trưởng có thể xem bracket, xem lịch trận, báo cáo kết quả và khiếu nại kết quả.

Các vai trò hỗ trợ gồm Trọng tài, Bình luận viên và Khán giả. Trọng tài theo dõi trận đấu và xử lý khiếu nại; Bình luận viên có thể nhận quyền theo dõi trận để phục vụ bình luận; Khán giả xem bracket, bảng xếp hạng và theo dõi kênh livestream của trận đấu.

\subsection{Quy trình nghiệp vụ}
\label{subsection:2.2.8}

Bên cạnh các biểu đồ use case, hệ thống còn có một số quy trình nghiệp vụ quan trọng kết hợp nhiều chức năng khác nhau. Trong phạm vi đồ án, ba quy trình tiêu biểu được lựa chọn để mô tả bằng biểu đồ hoạt động gồm quy trình quản lý và tham gia máy chủ, quy trình tổ chức và vận hành giải đấu, và quy trình đăng tải game.

\subsubsection{Quy trình nghiệp vụ Quản lý và Tham gia máy chủ}
\label{subsubsection:2.2.8.1}

\begin{figure}[H]
    \centering
    \includegraphics[width=\textwidth]{Hinhve/Chuong2/Activity/Quy-trinh-may-chu.png}
    \caption{Quy trình nghiệp vụ Quản lý và Tham gia máy chủ}
    \label{fig:2-8-activity-quy-trinh-may-chu}
\end{figure}

Hình \ref{fig:2-8-activity-quy-trinh-may-chu} mô tả quy trình từ khi Chủ máy chủ tạo máy chủ đến khi Người dùng có thể tham gia và truy cập các kênh được phép. Sau khi máy chủ được tạo, hệ thống khởi tạo các vai trò và kênh mặc định. Chủ máy chủ tiếp tục cấu hình thông tin, tạo kênh và thiết lập phân quyền cho máy chủ.

Người dùng có thể xem danh sách máy chủ và tham gia trực tiếp nếu máy chủ công khai. Với máy chủ yêu cầu xét duyệt, Người dùng gửi yêu cầu tham gia, hệ thống lưu yêu cầu ở trạng thái chờ và người có quyền duyệt sẽ chấp nhận hoặc từ chối. Khi được chấp nhận, hệ thống thêm Người dùng vào danh sách thành viên, gán vai trò mặc định và áp dụng quyền truy cập theo cơ chế phân quyền của máy chủ.

\subsubsection{Quy trình nghiệp vụ Tổ chức và Vận hành giải đấu}
\label{subsubsection:2.2.8.2}

\begin{figure}[H]
    \centering
    \includegraphics[width=\textwidth]{Hinhve/Chuong2/Activity/Quy-trinh-giai-dau.png}
    \caption{Quy trình nghiệp vụ Tổ chức và Vận hành giải đấu}
    \label{fig:2-9-activity-quy-trinh-giai-dau}
\end{figure}

Hình \ref{fig:2-9-activity-quy-trinh-giai-dau} thể hiện quy trình tổ chức một giải đấu trong máy chủ. Người tổ chức tạo giải đấu, cập nhật thông tin và mở đăng ký. Người tham gia đăng ký trực tiếp nếu là giải đấu cá nhân; với giải đấu theo đội, Đội trưởng tạo đội và thêm các thành viên cần thiết. Sau giai đoạn đăng ký, Người tổ chức chuyển giải đấu sang check-in để xác nhận các người chơi hoặc đội tham gia.

Khi giải đấu bắt đầu, hệ thống sinh bracket và danh sách trận đấu. Với mỗi trận, Người tổ chức bắt đầu trận đấu, hệ thống tạo workspace trận và phân quyền kênh cho người chơi, trọng tài hoặc các vai trò liên quan. Sau khi trận kết thúc, người tham gia báo cáo kết quả; nếu có khiếu nại, Trọng tài hoặc Quản trị giải đấu xem xét và ghi đè kết quả khi cần. Hệ thống cập nhật bracket, bảng xếp hạng và lặp lại quy trình cho đến khi toàn bộ trận đấu hoàn thành.

\subsubsection{Quy trình nghiệp vụ Đăng tải game}
\label{subsubsection:2.2.8.3}

\begin{figure}[H]
    \centering
    \includegraphics[width=\textwidth]{Hinhve/Chuong2/Activity/Quy-trinh-upload-game.png}
    \caption{Quy trình nghiệp vụ Đăng tải game}
    \label{fig:2-10-activity-quy-trinh-upload-game}
\end{figure}

Hình \ref{fig:2-10-activity-quy-trinh-upload-game} mô tả quy trình Nhà phát triển đưa game lên hệ thống. Trước hết, Nhà phát triển tạo thông tin game, nhập mô tả, hình ảnh và các metadata cần thiết. Hệ thống chuẩn hóa slug và tạo bản ghi game ở trạng thái nháp để Nhà phát triển tiếp tục tải lên bản build.

Sau khi Nhà phát triển tải lên file build dạng zip, hệ thống kiểm tra quyền sở hữu game, định dạng file, cấu trúc bundle và file \texttt{index.html}. Nếu build không hợp lệ, hệ thống trả lỗi để Nhà phát triển chỉnh sửa và tải lại. Nếu build hợp lệ, hệ thống lưu bundle, tạo đường dẫn chơi game, đánh dấu build sẵn sàng và cho phép gửi game vào hàng đợi kiểm duyệt. Quản trị viên kiểm tra nội dung game; nếu được duyệt, hệ thống cập nhật trạng thái phát hành và hiển thị game trên chợ game, ngược lại game được đưa về trạng thái cần chỉnh sửa.
